\section{Perturbations and identities on prescribed trees}
\label{pert:section}

Throughout this section the interaction and field are bounded, and the
depth $L\geq2$ is fixed.  The arguments apply both to the $N$-spin system
and to the reservoir in \cref{cavity:increment}. Its physical disorder
consists of independent bounded interactions with a Poisson count of
mean $O(N)$, together with $N$ independent bounded fields.
All expectations in a statement at a specified parameter vector are
taken with that vector held fixed.

We choose perturbations whose derivatives concentrate, while selecting
cavity increments at most an arbitrarily small error above the
unperturbed liminf. A Poisson insertion
then turns concentration into Taylor-coefficient identities. Independent
signs along the inserted paths select the prescribed splitting depths.
The use of concentration to obtain overlap identities has antecedents
in the self-averaging argument of Ghirlanda--Guerra~\cite{GG98}.
Related antecedents are the stochastic-stability relations of
Aizenman--Contucci~\cite{AC98} and their extensions to diluted
multioverlaps by Barra and De Sanctis~\cite{BarraDeSanctis07}.
Panchenko's spin-distribution identities~\cite{Pan13} provide the cavity
context. The bounded perturbation family and the identities on colored
finite trees needed here are constructed and proved below.

\subsection{The perturbation family}

A \emph{mark specification} consists of a finite number $n\geq0$ of
labeled Rademacher coordinates, a depth in $\{1,\ldots,L-1\}$ assigned
to each coordinate, two functions
$g,g_0:\{-1,1\}^n\to\mathbb Q\cap[-1,1]$, and a choice $S_0=1$ or
$S_0=\sigma_i$. Each coordinate is sampled at its assigned depth along
an auxiliary path. Distinct coordinates are independent; the value of
a coordinate is shared at a shared vertex and is independent at
different vertices. The case $n=0$ includes constant functions.
We use all such specifications, a countable collection because each
sign domain is finite. Products of signs and their rational averages
belong to this collection; these will select branching patterns.

The two parameters have different roles: differentiation in $u$ produces
a concentrated score, while the values of $t$ probe an insertion's Taylor
expansion. Taking $u$ smaller than every fixed power of $t$ will force
each limiting Taylor coefficient to vanish. For each specification
include types indexed by all integers $j\geq4$.
The type $\nu$ corresponding to this specification and $j$ has
\begin{equation}
 t_\nu=j^{-1},\qquad
 I_\nu=[j^{-j},2j^{-j}],\qquad u_\nu\in I_\nu.
 \label{pert:parameters}
\end{equation}
Enumerate these types by $\nu\in\mathbb N$ and choose constants
$c_\nu>0$ with $\sum_\nu c_\nu=1$.  Independently for each type, put
$n_\nu\sim\operatorname{Poisson}(c_\nu s_N)$, where $s_N=N^{3/4}$.
Each term has an independent uniform site $i\in[N]$ and an independent
copy of its mark process.  At a leaf it multiplies the spin weight by
\begin{equation}
 \psi_\nu=1+t_\nu g\sigma_i+u_\nu g_0S_0.
 \label{pert:factor}
\end{equation}
In particular $1/2<\psi_\nu<3/2$.  Counts, types, and sites belong to the
root; the mark variables are sampled with their product priors at their
assigned auxiliary depths.  Since the total count is
$\operatorname{Poisson}(s_N)$, this construction uses only finitely many
terms almost surely.  The factor bound permits termwise comparison
through the power means, as in \eqref{cavity:negligible}.

Write $\pi$ for the product of the uniform probability measures on the
intervals $I_\nu$, and $\mathbf u=(u_\nu)_\nu$. We use $\pi$ to select
deterministic parameter vectors, not as part of their quenched law.
Each interval has fixed positive length for a fixed type; concentration
estimates may therefore depend on its length.

\subsection{Concentration of each perturbation score}

For one type $\nu$, define on a tilted path, including its terminal
Gibbs configuration,
\begin{equation}
 D_\nu=\sum_{q=1}^{n_\nu}
 \frac{g_{0,q}S_{0,q}}
 {1+t_\nu g_q\sigma_{i_q}+u_\nu g_{0,q}S_{0,q}}.
 \label{pert:score}
\end{equation}
Let $\langle\cdot\rangle$ denote expectation along one tilted path
conditional on the root, and let $\mathbb E_{\mathbf u}$ average the
root at fixed $\mathbf u$.  Define
\begin{equation}
 e_{\nu,N}(\mathbf u)
 =\frac{1}{c_\nu s_N}
 \mathbb E_{\mathbf u}\Big\langle
 \big|D_\nu-\mathbb E_{\mathbf u}\langle D_\nu\rangle\big|
 \Big\rangle.
 \label{pert:error}
\end{equation}

\begin{lemma}[Averaged score concentration]
\label{pert:concentration}
For each fixed $L$ and each fixed type $\nu$,
\begin{equation}
 \int e_{\nu,N}(\mathbf u)\,\pi(d\mathbf u)\longrightarrow0.
 \label{pert:concentration-eq}
\end{equation}
The assertion holds for the reservoir as well as for the original
system.  Constants need not be uniform in $L$ or $\nu$.
\end{lemma}

\begin{proof}
Hold the other parameters fixed, write $u=u_\nu$, $n=n_\nu$, and put
$\lambda_\nu=c_\nu s_N$.  Choose an open interval $J_\nu$ containing
$I_\nu$ such that $|t_\nu|+|u|<1/2$ throughout its closure.  On this
interval, pathwise,
\begin{equation}
 |D_\nu|\leq2n,\qquad
 -4n\leq\partial_uD_\nu
 =-\sum_{q=1}^{n}
 \frac{(g_{0,q}S_{0,q})^2}{\psi_{\nu,q}^{2}}
 \leq0.
 \label{pert:score-bounds}
\end{equation}

\paragraph{Fluctuations along a tilted path.}
Let $X_0(u)$ be the random root log partition function before taking the
root expectation.  To include the last step in the differentiation,
regard the logarithm of the spin weight as $X_L$ and the spin partition
function as the step with exponent $m_L=1$.  Differentiation of one
backward step gives
\[
 \partial_uX_{d-1}=K_d(\partial_uX_d),\qquad
 \partial_u^2X_{d-1}
 =K_d(\partial_u^2X_d)
   +m_d\operatorname{Var}_{K_d}(\partial_uX_d).
\]
Thus, for $M_d=\langle D_\nu\mid\mathscr F_d\rangle$,
\begin{align}
 X_0'(u)&=\langle D_\nu\rangle,\label{pert:first-derivative}\\
 X_0''(u)&=\langle\partial_uD_\nu\rangle
   +\sum_{d=1}^{L}m_d
        \langle(M_d-M_{d-1})^2\rangle \nonumber\\
 &\geq m_1\operatorname{Var}(D_\nu\mid\mathscr F_0)-4n.
 \label{pert:second-derivative}
\end{align}
Here the conditional expectations and variances use the tilted path
law at the current $u$.  At fixed root all derivatives are justified by
the bounded score and its derivatives; the finite recursion then
justifies the displayed differentiation formulas.  Their root
expectations are justified by the finite moments of the Poisson counts.

Integrating \eqref{pert:second-derivative} over $I_\nu=[a,b]$ and using
$|X_0'|\leq2n$ yields
\begin{equation}
 \int_a^b\mathbb E_{\mathbf u}
       \operatorname{Var}(D_\nu\mid\mathscr F_0)\,du
 \leq\frac{4\lambda_\nu(1+b-a)}{m_1}.
 \label{pert:thermal-variance}
\end{equation}
Consequently the interval-averaged thermal mean absolute deviation,
divided by $\lambda_\nu$, is at most
\begin{equation}
 \left(\frac{4(1+|I_\nu|)}
 {m_1|I_\nu|\lambda_\nu}\right)^{1/2}.
 \label{pert:thermal-bound}
\end{equation}

\paragraph{Fluctuations of the root mean.}
It remains to concentrate $X_0'$ over the root. By
\eqref{pert:second-derivative}, the random function
\[
 G(u)=X_0(u)+2nu^2
\]
is convex on $J_\nu$.  Uniformly for $u\in J_\nu$ and all the other
parameters,
\begin{equation}
 \operatorname{Var}(G(u))\leq C_\nu N,
 \qquad |\mathbb E_{\mathbf u}G'(u)|\leq C_\nu\lambda_\nu.
 \label{pert:root-variance}
\end{equation}
To prove the first estimate, condition on the physical and total
perturbation counts. Conditional on the latter count, its types and sites
are independent labels; changing a label replaces a single bounded
factor. Apply the Efron--Stein variance inequality~\cite{EfronStein}
$\Var Z\le\frac12\sum_j\E(Z-Z^{(j)})^2$, where $Z^{(j)}$ is obtained
by independently resampling the $j$th input. For completeness, with
independent inputs $X_1,\ldots,X_k$, let $E_j$ integrate only $X_j$ and
let $\mathscr G_j=\sigma(X_1,\ldots,X_j)$. The Doob differences are
$D_j=\E[Z-E_jZ\mid\mathscr G_j]$. Orthogonality and Jensen give
\[
 \Var Z=\sum_j\E D_j^2
 \le\sum_j\E(Z-E_jZ)^2
 =\frac12\sum_j\E(Z-Z^{(j)})^2.
\]
This independent-input decomposition requires neither identical laws
nor a symmetric statistic; see Efron--Stein~\cite[Section~2]{EfronStein}.
Resampling one interaction, field, or
perturbation term changes $X_0$ by a bounded amount: this is immediate
for the spin weights and is preserved by each power mean.  When a term
is removed, its unused independent mark variables can be integrated
out, so the same assertion holds for addition or deletion of a term.
Conditional variances are therefore bounded by a constant times the
number of root variables.  Conditional means are Lipschitz functions
of the counts, with bounded Lipschitz constants.  If $Q,Q'$ are
independent copies of a Poisson variable and $b$ is $C$-Lipschitz, then
\[
 \operatorname{Var}(b(Q))
 =\tfrac12\mathbb E(b(Q)-b(Q'))^2
 \leq C^2\operatorname{Var}(Q).
\]
Applying this to the independent counts completes the $O(N)$ estimate.
The correction $2nu^2$ obeys the same estimates.  The derivative bound
in \eqref{pert:root-variance} follows directly from
\eqref{pert:score-bounds}.

\paragraph{Passing from a convex function to its derivative.}
Put
$\bar G=\mathbb E_{\mathbf u}G$ and $Z=G-\bar G$.  If $u\pm\rho$ lie
in $J_\nu$, forward and backward difference quotients imply
\begin{align*}
 |G'(u)-\bar G'(u)|
 &\leq\frac{|Z(u-\rho)|+2|Z(u)|+|Z(u+\rho)|}{\rho}\\
 &\quad+\bar G'(u+\rho)-\bar G'(u-\rho).
\end{align*}
The monotonicity of $\bar G'$ gives
\[
 \int_a^b[\bar G'(u+\rho)-\bar G'(u-\rho)]\,du
 \leq2\rho[\bar G'(b+\rho)-\bar G'(a-\rho)].
\]
Together with \eqref{pert:root-variance}, this proves
\begin{equation}
 \frac1{|I_\nu|}\int_{I_\nu}
 \mathbb E_{\mathbf u}|G'(u)-\mathbb E_{\mathbf u}G'(u)|\,du
 \leq C_\nu\left(\frac{\sqrt N}{\rho}
              +\frac{\rho\lambda_\nu}{|I_\nu|}\right).
 \label{pert:convex-bound}
\end{equation}
For all sufficiently large $N$, take $\rho=N^{-1/8}$, which lies
within the fixed neighborhood margin.  After division by
$\lambda_\nu=c_\nu N^{3/4}$, both terms tend to zero.  Removing the
correction to $G'$ costs at most
$C_\nu\mathbb E|n-\lambda_\nu|/\lambda_\nu
\leq C_\nu\lambda_\nu^{-1/2}$.  Thus the root contribution to
\eqref{pert:error} tends to zero.  Combining it with
\eqref{pert:thermal-bound} proves the result after integration over
$u_\nu$.  All estimates are uniform in the other parameters, so their
integration against the remaining product measure proves
\eqref{pert:concentration-eq}.
\end{proof}

\subsection{One sequence for low increments and all identities}

The parameter choices must satisfy concentration and a low cavity
increment simultaneously. Let $\underline F=\liminf_NF_N$ over all
system sizes for the bounded model. Recall that
$P_N^{L,\mathbf u}$ denotes the root-averaged perturbed log partition
function, and put
$\Delta_N(\mathbf u)=P_{N+1}^{L,\mathbf u}-P_N^{L,\mathbf u}$.

\begin{lemma}[Deterministic parameter selection]
\label[lemma]{pert:selection}
For every fixed $L\geq2$ and $\varepsilon>0$, there are sizes
$N_j\uparrow\infty$ and deterministic vectors $\mathbf u_j$ such that
\begin{enumerate}[label=\textup{(\roman*)}]
 \item $\Delta_{N_j}(\mathbf u_j)\leq \underline F+2\varepsilon$;
 \item for every fixed type $\nu$,
 $e_{\nu,N_j}(\mathbf u_j)\longrightarrow0$ in the reservoir;
 \item in the reservoir, the expectations of every finite product of multioverlaps on
 every fixed finite prescribed tree of depth $L$ converge along this
 sequence.
\end{enumerate}
\end{lemma}

\begin{proof}
\emph{Find low averaged increments.}
The uniform perturbation bound \eqref{cavity:negligible} gives
\[
 \liminf_N\frac1N\int P_N^{L,\mathbf u}\,\pi(d\mathbf u)=\underline F.
\]
Hence infinitely many $N$ satisfy
$\int\Delta_N\,d\pi\leq \underline F+\varepsilon$: otherwise telescoping the
averaged increments would contradict this liminf.  \Cref{cavity:increment}
supplies a constant $C$ such that
$|\Delta_N(\mathbf u)|\leq C$ uniformly.

\paragraph{Choose parameters with small score errors.}
For a finite collection of types, Lemma~\ref{pert:concentration} and
Markov's inequality show that the set where all its errors are small
has parameter measure tending to one.  More explicitly, choose
$N_j$ successively among the preceding infinite set so large that
\[
 \sum_{\nu\leq j}\int e_{\nu,N_j}(\mathbf u)\,\pi(d\mathbf u)
 \leq j^{-3}.
\]
Then the set $G_j$ where $e_{\nu,N_j}\leq j^{-1}$ for all $\nu\leq j$
has complement of measure $q_j\leq j^{-2}$.  Its conditional mean
increment is at most
\[
 \frac{\underline F+\varepsilon+Cq_j}{1-q_j}
 \leq \underline F+2\varepsilon
\]
for all sufficiently large $j$.  A vector in $G_j$ with increment at
most $\underline F+2\varepsilon$ can therefore be chosen. Relabeling the tail
proves (i) and (ii).

\paragraph{Make all finite-tree moments converge.}
At fixed $L$ the finite labeled tree
shapes and finite products of their multioverlaps form a countable
collection.  Each such product is bounded by one in absolute value.
A diagonal subsequence gives (iii), without affecting (i) or (ii).
\end{proof}

From now on $L$ is fixed and $N\to\infty$ means convergence along a
sequence from Lemma~\ref{pert:selection}. Expectations below are evaluated at
the selected deterministic vector, whose dependence on $N$ is
suppressed.  No uniformity in $L$, in the number of paths, or in a
perturbation type is asserted for these limits.

\subsection{Poisson insertion and analytic extraction}

Fix an old labeled tree $\mathcal T$, a distinguished leaf $a_*$, and
a bounded polynomial $f$ in multioverlaps of its old replicas.  The
polynomial and its coefficients are fixed as $N$ and the insertion
parameters vary. The
root-averaged expectation on this tree is denoted simply by
$\mathbb E f$.  The anchor marginal is always one tilted path.
Therefore \eqref{pert:error} implies, for each fixed type $\nu$,
\begin{equation}
 \left|\frac{\mathbb E[fD_\nu^{a_*}]
       -\mathbb E f\,\mathbb E\langle D_\nu\rangle}
       {c_\nu s_N}\right|
 \leq\|f\|_\infty e_{\nu,N}\longrightarrow0.
 \label{pert:score-identity}
\end{equation}

For a fixed mark specification, introduce an independent new uniform
site and an independent copy of its mark process, and set
$A=1+tg\sigma_i+ug_0S_0$.  Let $Y$ be its insertion recursion from
Lemma~\ref{tree:insertion}, and let $W_{\mathcal T}(A)$ be the updated-tree
density \eqref{tree:density}. The expectations in the following display use
the unmodified reservoir, extended by the independent insertion
variables with their prior kernels:
\begin{align}
 J_{N,\mathcal T}(t,u)
 &=\mathbb E\left[
 f\,\frac{g_0^{a_*}S_0^{a_*}}{A^{a_*}}
          W_{\mathcal T}(A)\right],\label{pert:palm-expression}\\
 H_N(t,u)
 &=J_{N,\mathcal T}(t,u)
   -\mathbb E f\,J_{N,\{a_*\}}(t,u).
 \label{pert:analytic-difference}
\end{align}
In the single-anchor expression $J_{N,\{a_*\}}$, the test is $1$.

\begin{lemma}[Palm identity]
\label{pert:palm}
If $\nu$ is a type with the specified marks, then
\begin{equation}
 |H_N(t_\nu,u_\nu)|\leq\|f\|_\infty e_{\nu,N}.
 \label{pert:palm-bound}
\end{equation}
For a fixed mark specification the function $H_N(t,u)$ is the same
for all its types $j$: its base reservoir contains the full original
Poisson families, including all types.
\end{lemma}

\begin{proof}
Apply the Poisson add-one identity to the sum over type-$\nu$ terms
in $D_\nu$.  After division by its mean $c_\nu s_N$, the remaining
count is again Poisson with mean $c_\nu s_N$, and all the other counts
are unchanged.  Thus the remaining model has exactly the original
reservoir law.  The distinguished term becomes an independent new
insertion with parameters $t_\nu,u_\nu$.  Its score is
$g_0^{a_*}S_0^{a_*}/A^{a_*}$.  The updated tree density is
\eqref{tree:density}, so the anchor factor $A^{a_*}$ cancels.
This proves
\[
 \frac{\mathbb E[fD_\nu^{a_*}]}{c_\nu s_N}
 =J_{N,\mathcal T}(t_\nu,u_\nu),\qquad
 \frac{\mathbb E\langle D_\nu\rangle}{c_\nu s_N}
 =J_{N,\{a_*\}}(t_\nu,u_\nu).
\]
Equation~\eqref{pert:palm-bound} now follows from
\eqref{pert:score-identity}.  In particular no type has been deleted
from the law of the base reservoir.  This proves the last assertion.
\end{proof}

\begin{lemma}[Vanishing insertion coefficients]
\label[lemma]{pert:analytic}
For every mark specification included in the perturbation family, old tree, anchor, and bounded
overlap-polynomial test $f$, every Taylor coefficient at $t=0$ of
$H_N(t,0)$ tends to zero along the selected sequence.
\end{lemma}

\begin{proof}
Lemma~\ref{tree:extension-expansion} gives a complex neighborhood
$|t|,|u|<\rho$ on which the functions in
\eqref{pert:analytic-difference} are analytic and uniformly bounded
in $N$ and in the selected parameter vector.  Shrinking $\rho$ gives
a uniform bound on their $u$-derivatives as well.  These constants may
depend on the fixed tree, $L$, and $\|f\|_\infty$, but not on the type
index $j$ in the mark specification.

Each Taylor coefficient of $H_N(t,0)$ is a finite linear combination
of multioverlap moments and products of two such moments on prescribed trees.
Their coefficients are determined by the fixed tree, exponents, mark law,
and polynomial test, and are independent of $N$. Indeed the canceled
anchor expansion supplies finitely many extensions at any fixed
order; averaging the independent insertion site turns its spin
product into a multioverlap.  These coefficients converge by
Lemma~\ref{pert:selection}(iii).  The uniform analytic bound then
implies, by the geometric bound on the Taylor tails, that
$H_N(t,0)$ converges uniformly on each smaller disk to an analytic
function $H(t)$.

For every sufficiently large fixed $j$, the corresponding type and
Lemma~\ref{pert:palm} give
\[
 \limsup_{N\to\infty}|H_N(j^{-1},0)|
 \leq C\limsup_{N\to\infty}|u_{\nu,N}|
 \leq2Cj^{-j}.
\]
Here the score error tends to zero with $j$ fixed, and the same
function $H_N$ is used for every $j$, by Lemma~\ref{pert:palm}.
Hence $|H(j^{-1})|\leq2Cj^{-j}$.  If the first nonzero Taylor
coefficient of $H$ had degree $q$, then
$|H(j^{-1})|\geq c j^{-q}$ for all sufficiently large $j$, a
contradiction.  Thus $H$ vanishes identically near zero, proving the
claim coefficient by coefficient.
\end{proof}

\subsection{Prescribing the colored tree}

The anchor is called a \emph{spin anchor} when $S_0=\sigma_i$, and a
\emph{marker anchor} when $S_0=1$.  Added colors are ordered
$1,\ldots,k$; the anchor has color $0$.  They occupy distinct replica
positions, none equal to the anchor, according to the canceled-anchor
extension rule of Lemma~\ref{tree:extension-expansion}.  A constraint
$\mathcal C$ prescribes any consistent collection of their pairwise
split depths, involving only the added colors and the anchor.
It imposes no extra restrictions involving the other old replicas.
Let $\mathfrak E(\mathcal T,a_*,\mathcal C)$ be the allowed extension
histories, with coefficients $c_E$ from that lemma.  For a history
$E$, set
\[
 R_E=\frac1N\sum_{i=1}^N
 \begin{cases}
 \displaystyle\sigma_i^{a_*}\prod_{j=1}^k\sigma_i^{a_j(E)},
       &\text{spin anchor},\\[2pt]
 \displaystyle\prod_{j=1}^k\sigma_i^{a_j(E)},
       &\text{marker anchor}.
 \end{cases}
\]
For $k=0$ these are respectively $R_{\{a_*\}}$ and $1$.

\begin{theorem}[Identity on prescribed trees]
\label{pert:tree-identity}
\label{pert:identity}
For the perturbation family defined above, along every sequence from
Lemma~\ref{pert:selection},
\begin{equation}
 \sum_{E\in\mathfrak E(\mathcal T,a_*,\mathcal C)}
       c_E\,\mathbb E[fR_E]
 -\mathbb E f
  \sum_{E\in\mathfrak E(\{a_*\},a_*,\mathcal C)}
       c_E\,\mathbb E[R_E]
 \longrightarrow0.
 \label{pert:tree-identity-eq}
\end{equation}
This holds for every fixed old tree, either kind of anchor, every
bounded polynomial $f$ in its multioverlaps, every finite number of
ordered added colors, and every prescribed constraint $\mathcal C$.
The empty addition is allowed.  The error is taken to zero at fixed
$L$, fixed tree, and fixed constraint, before any limit in $L$.
\end{theorem}

\begin{proof}
For $k=0$, the claim is the vanishing constant coefficient from
\cref{pert:analytic}, using constant marks. For $k\ge1$, we explain
how to extract distinct colored marks. Fix rational-valued functions
$g_0,g_1,\ldots,g_k$, all bounded by one and constructed jointly from
the same finite collection of auxiliary marks.  The mixed insertion
uses $A=1+\sum_{j=1}^kz_jg_j\sigma_i$.  Its mixed derivative in all
$z_j$ is the symmetric multilinear form obtained by polarizing the
degree-$k$ coefficient for a single $g$.  More explicitly, if
$Q_N(g)=k![t^k]H_N(t,0)$, then the mixed derivative is
\[
 \frac{k^k}{2^k k!}
 \sum_{\epsilon\in\{-1,1\}^k}
    \left(\prod_{j=1}^k\epsilon_j\right)
 Q_N\left(k^{-1}\sum_{j=1}^k\epsilon_jg_j\right).
\]
Every argument of $Q_N$ is rational-valued and bounded by one, so each
resulting pair $(g,g_0)$ belongs to the fixed mark family.  Lemma
\ref{pert:analytic} makes the displayed mixed derivative tend to zero.
The mixed derivative expansion is precisely the ordered extension
rule: there is no additional factorial in its coefficients.

It remains to arrange that the expectation over the marks selects
exactly the extensions satisfying $\mathcal C$.  For a pair of colors
$a,b\in\{0,\ldots,k\}$ and a depth $1\leq d\leq L-1$, the event
that their split depth is at least $d$ is the event that they share
their vertex at depth $d$.  Introduce a new independent Rademacher
coordinate at that depth and give both $g_a$ and $g_b$ its factor.
Its contribution to their joint mark expectation is $1$ if the
vertex is shared and $0$ otherwise.  For a finite collection of such
requirements, use different independent coordinates for the
different requirements and define each $g_a$ as the product of its
assigned factors.  Then
\begin{equation}
 \mathbb E_{\rm marks}
       \left[g_0^{a_*}\prod_{j=1}^kg_j^{a_j(E)}\right]
 =\prod_{(a,b,d)\ {\rm required}}
       \mathbf1_{\{a,b\text{ share their depth-}d\text{ vertex}\}}.
 \label{pert:mark-selector}
\end{equation}
All marks in this calculation are independent of the base reservoir.

Let $H_d(a,b)$ denote the sharing indicator in
\eqref{pert:mark-selector}, with $H_0=1$ and $H_L=0$ for distinct
replica positions. These are positions in the sampled tree, even when
their spin configurations happen to agree. Exact splitting at depth $d$ is expressed by
$H_d-H_{d+1}$, for $0\leq d\leq L-1$.  Products and finite linear
combinations therefore select any prescribed consistent colored
shape or collection of allowed shapes. For each term of this finite
linear combination, the functions $g_0,\ldots,g_k$ are products of signs
and hence take values in $\{-1,1\}$. These functions and their
polarization averages belong to the mark family fixed at the start of
the section. Constant marks cover $k=0$.

Finally apply the mixed canceled-anchor expansion to the difference
\eqref{pert:analytic-difference}, average the marks using the
selector, and average the independent uniform insertion site.  The
latter average is exactly $R_E$.  Lemma~\ref{pert:analytic} and
polarization make the coefficient difference tend to zero, yielding
\eqref{pert:tree-identity-eq}.  Finite linear combinations extend the
conclusion from overlap monomials to the stated tests.  The
construction uses only marks through depth $L-1$: the extension rule
never reuses the anchor or an already used color, so no test of
coincidence at depth $L$ is required.
\end{proof}

\begin{remark}[Order of the choices]
The mark specifications and their positive weights $c_\nu$ are fixed
first, for a fixed $L$.  Parameter averaging then gives one selected
sequence on which all types concentrate and all finite overlap
moments converge.  The analytic argument takes $N\to\infty$ with
each $j$ fixed before sending $j\to\infty$.  Only after the resulting
prescribed-tree identities have been established is $L$ allowed to
increase.  No concentration estimate uniform over the countable
family is needed.
\end{remark}
